%% Problem 1 (AHP site scoring) and Problem 2 (facility location / covering
%% models): what each problem is, why it is needed, and how it is solved.
%% Build:  cd paper && pdflatex prob1_prob2 && pdflatex prob1_prob2
\documentclass[11pt]{article}

\usepackage[margin=1in]{geometry}
\usepackage{amsmath,amssymb}
\usepackage{booktabs}
\usepackage[hidelinks]{hyperref}
\usepackage{parskip}

\title{Problem 1 \& Problem 2: What, Why, How and Solution}
\author{Dynamic EMS Dispatch \& Station Placement Optimization}
\date{\today}

\begin{document}
\maketitle

% ===========================================================================
\section*{Problem 1 --- AHP Site Scoring}
% ===========================================================================

\subsection*{What}
A multi-criteria ranking of every candidate station site, using Saaty's
Analytic Hierarchy Process (AHP). Three benefit criteria are scored per
candidate zone: \textbf{Incident Density} (how many emergencies happen there),
\textbf{Equity Gap} (distance to the nearest station Seattle already
operates, farther read as higher unmet need), and \textbf{Coverage Potential}
(how much demand that site could reach within the response-time limit).
These are combined into one \emph{priority score} per zone.

\subsection*{Why}
The facility-location problem that follows (Problem 2) can have more than one
optimal station set of the same size: several different subsets of zones may
all achieve full coverage with the same minimum station count. Picking among
them arbitrarily throws away information the problem has available --- some
sites are simply more strategically important than others (busier, farther
from existing help, able to reach more people). AHP gives a transparent,
auditable way to rank sites on multiple competing factors at once, with a
built-in check that the ranking is not self-contradictory, rather than
weighting the criteria by guesswork.

\subsection*{How}
\begin{enumerate}
  \item Each team member independently rates the three criteria pairwise, on
    Saaty's 1--9 intensity scale, producing a $3\times3$ reciprocal matrix
    $A^{(k)}$ per member ($A^{(k)}_{ji} = 1/A^{(k)}_{ij}$, diagonal 1).
  \item The members' matrices are combined by element-wise \textbf{geometric
    mean} --- the only averaging rule that keeps the result reciprocal:
    \[
      \bar A_{ij} = \Big(\textstyle\prod_{k=1}^{n} A^{(k)}_{ij}\Big)^{1/n}
    \]
  \item The \textbf{priority weights} $w$ are the principal eigenvector of
    $\bar A$, normalised to sum to 1: $\bar A w = \lambda_{\max} w$.
  \item \textbf{Consistency} is checked before the weights are trusted:
    \[
      \mathrm{CI} = \frac{\lambda_{\max}-n}{n-1}, \qquad
      \mathrm{CR} = \frac{\mathrm{CI}}{\mathrm{RI}_n}, \qquad
      \text{accepted iff } \mathrm{CR} < 0.10
    \]
  \item Each candidate zone's three raw criterion values are
    \textbf{max-normalised} ($\hat x_{ic} = x_{ic}/\max_i x_{ic}$) and combined
    into a \textbf{priority score}:
    \[
      s_i = \sum_{c} w_c\, \hat x_{ic}
    \]
  \item The score orders the candidate list and enters Problem 2's objective
    only as an \emph{infinitesimal tie-break} ($\varepsilon < 1/n_{\text{candidates}}$)
    --- it can never justify opening an extra station; it only chooses which
    equally-good station set is reported.
\end{enumerate}

\subsection*{Solution}
Four team members rated the three criteria independently on the Saaty scale;
every individual matrix, and the aggregate, came out \emph{exactly}
consistent:

\begin{table}[h]
\centering
\small
\begin{tabular}{@{}lccccccc@{}}
\toprule
 & \multicolumn{3}{c}{\textbf{Saaty ratings}} & \multicolumn{3}{c}{\textbf{Weights (\%)}} & \\
\cmidrule(lr){2-4}\cmidrule(lr){5-7}
\textbf{Member} & D:E & D:C & E:C & D & E & C & \textbf{CR} \\
\midrule
Aadith Krishna       & 3    & 6    & 2    & 66.7 & 22.2 & 11.1 & 0.000 \\
Aditya Santosh       & 2    & 6    & 3    & 60.0 & 30.0 & 10.0 & 0.000 \\
Suryansh Ram Menon   & 4    & 8    & 2    & 72.7 & 18.2 &  9.1 & 0.000 \\
Arjungopal Anilkumar & 3    & 9    & 3    & 69.2 & 23.1 &  7.7 & 0.000 \\
\midrule
\textbf{Aggregate}   & 2.91 & 7.14 & 2.45 & \textbf{67.4} & \textbf{23.1} & \textbf{9.4} & 0.000 \\
\bottomrule
\end{tabular}
\end{table}
\noindent (D = Incident Density, E = Equity Gap, C = Coverage Potential;
$\lambda_{\max}=3.000$, $\mathrm{CI}=0.000$, $\mathrm{RI}_3=0.58$.)

The team therefore weights raw call volume (\textbf{67.4\%}) far above
distance from existing help (\textbf{23.1\%}), and ranks the ability to
blanket several zones from one site (\textbf{9.4\%}) least important of the
three. Applying these weights, the four highest-priority candidate zones are
\textbf{ZIP 98104} (score 0.778), \textbf{98144} (0.579), \textbf{98101}
(0.534) and \textbf{98122} (0.512).

A perfectly zero CR for every member is not a coincidence: in each member's
matrix the Density-vs-Coverage rating happens to equal the product of the
other two judgements, which is exactly the algebraic condition for perfect
consistency, so the check has no contradiction left to find on this input. An
independent randomised test (200 random reciprocal matrices, checked against
an independently computed eigenvector) confirms separately that the
consistency machinery itself is live and would catch a real contradiction.
This ranking is what Problem 2 below uses as its tie-break.

% ===========================================================================
\section*{Problem 2 --- Facility Location (Covering Models)}
% ===========================================================================

\subsection*{What}
The strategic decision of \emph{where to place ambulance stations}. Two
complementary integer programs are solved over the same candidate sites and
the same travel-time matrix:
\begin{itemize}
  \item \textbf{Mode A -- Set Covering (LSCP):} the fewest stations such that
    every demand zone is reached within the response-time limit.
  \item \textbf{Mode B -- Maximal Covering (MCLP):} with the station budget
    fixed at $P$, the placement that covers the most demand possible.
\end{itemize}

\subsection*{Why}
These two models answer two different real questions, and together they give
the full picture a single model cannot. LSCP answers \emph{``what is the
cheapest network that leaves nobody outside the limit?''} --- the right
question when full coverage is the goal and cost is secondary. MCLP answers
the complementary, budget-constrained question: \emph{``we can only afford
$P$ stations; where do they do the most good?''} --- the right question once
the budget, not the coverage promise, is the binding constraint (for
instance, because the marginal station buys rapidly diminishing extra
coverage, which only the MCLP curve reveals). Reporting a point estimate from
either alone would also hide how sensitive it is to the response-time limit
and to the assumed travel speed/circuity, so the models are re-solved across
a grid of both to report a range rather than a single number.

\subsection*{How}
\begin{enumerate}
  \item \textbf{Candidate sites and demand zones} are the same 28 material
    ZIP-code zones from the data pipeline, each represented by its
    demand-weighted centroid (the mean coordinate of real incidents in that
    zone).
  \item \textbf{Travel time} $t_{ij}$ between every candidate $i$ and zone $j$
    is estimated from the great-circle (Haversine) distance $d_{ij}$, inflated
    by a road circuity factor $\kappa$, divided by an effective speed $v$, plus
    a fixed turnout time $\tau$:
    \[
      t_{ij} = \tau + \frac{60\,\kappa\, d_{ij}}{v}
    \]
  \item A \textbf{coverage set} $N(j) = \{\,i : t_{ij} \le T\,\}$ is built for
    the chosen response-time limit $T$; any zone with $N(j)=\varnothing$ is
    flagged as structurally uncoverable and excluded from the constraints
    rather than left to make the whole program infeasible.
  \item \textbf{Mode A} solves
    \[
      \min_y \sum_{i} y_i \quad\text{s.t.}\quad \sum_{i\in N(j)} y_i \ge 1\ \ \forall j,
      \qquad y_i \in \{0,1\}
    \]
    (plus the Problem 1 tie-break term), returning the minimum station count
    and which sites to open.
  \item \textbf{Mode B} solves, for a chosen budget $P$,
    \[
      \max_{y,z} \sum_{j} d_j z_j \quad\text{s.t.}\quad
      z_j \le \sum_{i\in N(j)} y_i\ \ \forall j, \quad \sum_i y_i = P,
      \quad y_i,z_j\in\{0,1\}
    \]
    where $d_j$ is zone $j$'s demand, re-solved at every $P$ from 1 upward to
    trace the diminishing-returns curve.
  \item Both ILPs are solved exactly (HiGHS/CBC via PuLP), in milliseconds at
    this instance size (28 candidates $\times$ 28 zones), and are
    cross-checked against brute-force enumeration on synthetic small-city
    instances to confirm the solver reaches the true optimum.
  \item The whole Mode A solve is then repeated across a grid of $(v,\kappa)$
    values, since response times here are modelled rather than observed, so
    the reported station count is a range (e.g.\ seven to seventeen) rather
    than a single point estimate.
\end{enumerate}

\subsection*{Solution}
Over the study period (79{,}152 EMS incidents, 28 material demand zones),
Mode A was solved at each of three response-time limits and benchmarked
against Seattle's own 36-station network under identical travel assumptions:

\begin{table}[h]
\centering
\small
\begin{tabular}{@{}lrrrr@{}}
\toprule
\textbf{Limit} & \multicolumn{2}{c}{\textbf{Optimised}} & \multicolumn{2}{c}{\textbf{Seattle's real network}} \\
\cmidrule(lr){2-3}\cmidrule(lr){4-5}
 & Stations & Demand & Stations & Demand \\
\midrule
5 min  & 16 & 100.00\% & 36 & 99.72\% \\
8 min  &  7 & 100.00\% & 36 & 100.00\% \\
10 min &  5 & 100.00\% & 36 & 100.00\% \\
\bottomrule
\end{tabular}
\end{table}

At the headline \textbf{8-minute} limit, \textbf{7 stations} reach full
coverage of all 28 material zones -- the same 100\% coverage Seattle's real
36-station network achieves, from under a fifth of the sites. At the tighter
5-minute limit the optimised plan (16 stations) even edges out the real
network's coverage (100.00\% vs.\ 99.72\%) with fewer than half its stations.

\textbf{Mode B's} diminishing-returns curve justifies the budget-constrained
framing: the first station alone covers \textbf{51\%} of demand, five
stations cover \textbf{98\%}, and the remaining two (of seven) buy only the
last \textbf{1.8\%}.

\textbf{Robustness.} Because response times are modelled, not measured, Mode
A was re-solved over a $4\times4$ grid of speed (30--45 km/h) $\times$
circuity (1.3--1.6). The station count required for full coverage ranges
from \textbf{7 to 17} (a factor of 2.4) across that grid, while demand
covered stays at \textbf{100\%} in every one of the 16 cells. Seventeen is
still under half of Seattle's 36, so the \emph{directional} finding, fewer
sites than the real network, survives even the most pessimistic reading;
only the exact count is assumption-dependent.

\textbf{Why fewer stations suffice.} Three factors explain the gap, and none
of them means Seattle Fire is badly run: (1) the real 36 solve a broader
problem -- fire suppression, hazmat and rescue, not EMS alone, and two of the
36 are not even numbered stations (department HQ and Harborview Medical
Center); (2) the real network accreted historically over a century of
annexation and politics, while the optimiser starts from a blank map; and (3)
coverage here is evaluated at 28 demand-weighted centroids rather than at
each of the 79{,}152 individual calls, which is itself optimistic -- pricing
in that within-zone travel raises the 8-minute answer from 7 stations to 12
(quantified, not asserted, in the paper's Limitations section). The honest
reading is therefore not ``Seattle should close 29 fire stations,'' but that
\emph{EMS demand in Seattle is concentrated enough that a small,
demand-sited ambulance fleet can cover it} -- an argument about ambulance
placement, not fire-station count.

\end{document}
